Three-dimensional Gaussian Splatting (3DGS) is an explicit, point-based scene representation introduced by Kerbl, Kopanas, Leimkühler, and Drettakis at SIGGRAPH 2023 {[}Kerbl et al.~2023, \emph{ACM Transactions on Graphics} 42(4){]}. A scene is encoded as a set of millions of anisotropic 3D Gaussian primitives that are differentiably rasterized to images in real time. On the Mip-NeRF 360 benchmark of nine unbounded real scenes, vanilla 3DGS reaches PSNR \(27.21\) dB, SSIM \(0.815\), and LPIPS \(0.214\) at \(134\) FPS at \(1080p\) on an NVIDIA RTX A6000, while training in \(30\)--\(60\) minutes --- the first method to match Mip-NeRF 360 quality at real-time rates. This combination of photorealism and speed reshaped neural rendering and triggered an explosion of follow-up work that, as of early \(2026\), already exceeds the entire pre-3DGS NeRF literature in volume. This survey covers 3DGS work from August \(2023\) through early \(2026\), including its EWA-splatting and image-based-rendering antecedents. Methods are organized along five axes: primitive geometry (vanilla, 2DGS, surfels, convexes, tetrahedra), anti-aliasing (Mip-Splatting, Multi-Scale 3DGS, RadSplat), compression (CompGS, Compact 3DGS, EAGLES, Niedermayr), optimization paradigm (per-scene vs feed-forward --- pixelSplat, MVSplat), and downstream task (SLAM, dynamics, surfaces, generation). The lineage spans EWA splatting {[}Zwicker et al.~2001, 2002{]}, COLMAP image-based rendering {[}Schönberger and Frahm 2016{]}, and neural radiance fields {[}Mildenhal\ldots{}
